% ============================================================
\section{Rel (Relations)}
As construções na categoria de relações focam-se em combinar e intersetar conexões genéricas. O foco principal é garantir que a conectividade relacional seja transportada de forma consistente para as novas estruturas, mesmo na ausência de restrições algébricas adicionais.

\subsection{Fundamentação Teórica}
As principais construções universais implementadas para Rel são:
\begin{itemize}
    \item \textbf{Produtos:} O produto cartesiano $A \times B$ equipado com a relação onde $(a,b) \sim (a',b')$ se e só se $a \sim_A a'$ e $b \sim_B b'$.
    \item \textbf{Equalizadores:} Sub-relação correspondente aos elementos onde os dois morfismos paralelos são iguais.
\end{itemize}

\subsection{Implementação Computacional}
O produto tensorial (ou de Kronecker) aplica-se perfeitamente para combinar estas relações de forma vetorial em MATLAB.

\begin{lstlisting}[caption={Construção do Produto de Relações}]
function P = rel_product(MA, MB)
    % Produto relacional via produto de Kronecker
    P_matrix = kron(MA, MB);
    P.matrix = P_matrix;
    P.n_elements = size(P_matrix, 1);
end
\end{lstlisting}